% =====================================================================
%  FICHE 11 — THEME 3 — Exercices MOYEN — Équilibrer une équation redox
% =====================================================================
\documentclass[11pt,a4paper]{article}
\usepackage[utf8]{inputenc}
\usepackage[T1]{fontenc}
\usepackage{lmodern}
\usepackage[french]{babel}
\usepackage{amsmath,amssymb}
\usepackage[version=4]{mhchem}
\usepackage{siunitx}
\sisetup{output-decimal-marker={,},per-mode=symbol}
\usepackage{xcolor}
\usepackage{array}
\usepackage{booktabs}
\usepackage{enumitem}
\usepackage[most]{tcolorbox}
\usepackage[a4paper,margin=2.1cm,headheight=26pt,headsep=10pt]{geometry}
\usepackage{fancyhdr}
\usepackage[hidelinks]{hyperref}

\definecolor{primary}{RGB}{146,95,160}
\definecolor{primarydark}{RGB}{92,52,104}
\newcommand{\NO}{\ensuremath{\mathrm{N.O.}}}

\newcounter{exo}
\newtcolorbox{exo}[1][]{enhanced,breakable,boxrule=0.9pt,arc=2pt,colback=primary!4,
  colframe=primary,fonttitle=\bfseries,coltitle=white,left=3mm,right=3mm,top=2.2mm,bottom=2.2mm,
  title={Exercice~\refstepcounter{exo}\theexo\ifx\relax#1\relax\relax\else\ \textmd{--- \itshape #1}\fi}}

\pagestyle{fancy}\fancyhf{}
\fancyhead[L]{\small\textcolor{primarydark}{\textbf{Fiche 11 · Thème 3} — Équilibrer une équation redox}}
\fancyhead[R]{\small\textcolor{primarydark}{\textsc{Moyen}}}
\fancyfoot[C]{\small\thepage}
\renewcommand{\headrulewidth}{0.8pt}
\renewcommand{\headrule}{\hbox to\headwidth{\color{primary}\leaders\hrule height \headrulewidth\hfill}}

\begin{document}
\begin{center}
{\LARGE\bfseries\color{primarydark}Thème 3 — Niveau Moyen}\\[2pt]
{\small\itshape Combiner les demi-équations, gérer \ce{H+}/\ce{H3O+}/\ce{OH^-} et sommer les coefficients.}\\[-4pt]
{\color{primary}\rule{\textwidth}{1pt}}
\end{center}

\begin{exo}[assembler deux demi-équations — milieu acide]
On donne les deux demi-équations déjà équilibrées :
\[
\ce{NO3^- + 4 H+ + 3 e^- -> NO + 2 H2O} \qquad ; \qquad \ce{Fe^{2+} -> Fe^{3+} + e^-}.
\]
Combine-les pour obtenir l'\textbf{équation globale} (élimine les électrons), puis contrôle les
charges.
\end{exo}

\begin{exo}[équilibrer avec des ions hydronium — d'après annales 2024]
On considère la réaction d'oxydo-réduction non pondérée :
\[
\ce{PbO2 + Cr^{3+} + H2O -> Pb^{2+} + Cr2O7^{2-} + H3O+}.
\]
\begin{enumerate}[label=(\alph*),leftmargin=2em,topsep=1pt,itemsep=1pt]
  \item Identifie les deux couples et le nombre d'électrons de chaque demi-équation.
  \item Équilibre chaque demi-équation en milieu acide (avec \ce{H+} et \ce{H2O}), puis combine-les.
  \item Récris l'équation globale avec \ce{H3O+} (chaque \ce{H+} devient \ce{H3O+} moyennant une \ce{H2O}).
        Quel est alors le \textbf{coefficient de l'ion hydronium} \ce{H3O+} ?
\end{enumerate}
\end{exo}

\begin{exo}[une demi-équation en milieu basique]
Équilibre en \textbf{milieu basique} la demi-équation de réduction du couple \ce{MnO4^-}/\ce{MnO2} :
\[
\ce{MnO4^- -> MnO2}.
\]
Méthode : équilibre d'abord comme en milieu acide, puis neutralise les \ce{H+} par des \ce{OH^-}.
Contrôle les charges à la fin.
\end{exo}

\begin{exo}[permanganate et eau oxygénée — d'après annales 2022]
Détermine les coefficients stœchiométriques entiers de la réaction (milieu acide) :
\[
\ce{MnO4^- + H+ + H2O2 -> Mn^{2+} + H2O + O2}.
\]
Écris d'abord les deux demi-équations (\ce{MnO4^-}/\ce{Mn^{2+}} et \ce{O2}/\ce{H2O2}), puis combine.
\end{exo}

\begin{exo}[somme des coefficients — d'après livre QCM 13]
On donne la réaction non pondérée (milieu acide) :
\[
\ce{Sn + NO3^- + H+ -> SnO2 + NO2 + H2O}.
\]
\begin{enumerate}[label=(\alph*),leftmargin=2em,topsep=1pt,itemsep=1pt]
  \item Équilibre complètement cette équation.
  \item Calcule la \textbf{somme des coefficients des réactifs} et celle \textbf{des produits}. Laquelle
        est la plus grande ?
\end{enumerate}
\end{exo}

\end{document}
